\section{Related Work}
\label{sec:related}

% -----------------------------------------------------------------------------
\subsection{Biomedical Image Segmentation}
Biomedical image segmentation encompasses diverse clinical imaging modalities with distinct visual and structural properties~\cite{litjens2017survey}. Beginning with the foundational U-Net architecture~\cite{ronneberger2015unet} and its multi-scale variants such as UNet++~\cite{zhou2018unetplusplus} and ResUNet~\cite{zhang2018road}, deep convolutional neural networks have set the benchmark in medical image analysis. To capture ambiguous boundary transitions and multi-scale context, attention-guided models have emerged, including Attention U-Net~\cite{oktay2018attention}, PraNet~\cite{fan2020pranet} (which applies reverse attention for polyp boundary modeling), and HarDNet-MSEG~\cite{huang2021hardnet}. Recently, vision transformer architectures such as TransUNet~\cite{chen2021transunet}, Swin-Unet~\cite{cao2022swin}, and Polyp-PVT~\cite{dong2021polyp} have incorporated self-attention mechanisms to capture long-range contextual relationships.

Despite their expressive power, standard feedforward architectures remain inherently single-pass: intermediate representations cannot be re-examined or progressively modulated by downstream contextual hypotheses. In challenging clinical regimes---such as flat sessile polyps, thin capillary vessels, or overlapping cell nuclei---single-pass feedforward models frequently struggle to separate subtle lesion boundaries from healthy mucosal textures.

% -----------------------------------------------------------------------------
\subsection{Feedback Mechanisms and Recurrent Vision Networks}
In biological vision, neuroanatomical feedback pathways connecting higher cortical visual areas to early visual processing stages substantially outnumber feedforward projections~\cite{felleman1991distributed, gilbert2013top}. Inspired by this biological architecture, recurrent feedback networks simulate iterative visual cognition in deep neural networks~\cite{liang2015recurrent, zamir2017feedback}. In general computer vision, Feedback Networks~\cite{zamir2017feedback} and recurrent convolutional networks (RCNN)~\cite{pinheiro2014recurrent} have shown that routing higher-level predictions to earlier feature extractors significantly aids spatial disambiguation.

In medical image segmentation, Recurrent U-Net~\cite{alom2018recurrent}, BCDU-Net~\cite{azad2019bi}, and the Feature Attention Network (FANet)~\cite{tomar2021fanet} introduced iterative feedback loops across training cycles. FANet, in particular, proposed the \texttt{MixPool} module, which merges the previous epoch's binary prediction mask with intermediate encoder and decoder features to guide spatial attention.

\textbf{The Overlooked Research Gap:}
Existing recurrent segmentation literature operates under the unverified heuristic assumption that re-injecting prediction masks into intermediate layers automatically facilitates iterative boundary refinement. These studies treat recurrent coupling as a black box, overlooking the underlying gradient dynamics. In this work, we demonstrate that hard-thresholded mask injection creates a catastrophic \textit{Feedback Trap}: it produces vanishing gradients in internal attention branches, allows unchecked recursive backpropagation to distort encoder Batch Normalization statistics, and locks the network into runaway false-positive hallucinations.

% -----------------------------------------------------------------------------
\subsection{Loss Functions and Class Imbalance in Medical Vision}
Segmentation of subtle lesions, micro-vessels, and small nuclei is heavily impacted by extreme class imbalance, where background pixels vastly outnumber foreground targets~\cite{salehi2017tversky}. While standard Cross-Entropy loss is easily biased toward the background class, overlap-based losses such as the Dice loss~\cite{milletari2016vnet} and soft Jaccard loss~\cite{rahman2016optimizing} mitigate imbalance by optimizing spatial overlap directly.

To penalize boundary errors and suppress false positives, asymmetric formulations have been introduced. The Tversky loss~\cite{salehi2017tversky} generalizes the Dice index with tunable weights $\alpha$ and $\beta$ to penalize false positives and false negatives independently. The Focal Tversky loss~\cite{abraham2019novel} modulates hard examples, while the Unified Focal loss~\cite{yeung2022unified} balances precision and recall. 

\textbf{Our Mechanistic Finding:}
Prior research has attempted to resolve false-positive over-segmentation primarily by designing increasingly complex loss functions. However, our factorial experiments reveal that such loss engineering is rendered largely ineffective when paired with conventional recurrent feedback loops. By resolving the structural root cause of the Feedback Trap through gradient decoupling, our proposed GDF-Net unleashes the corrective power of asymmetric boundary optimization, ensuring stable convergence across diverse biomedical imaging domains.

% -----------------------------------------------------------------------------
\subsection{Gradient Decoupling and Stop-Gradient Dynamics}
In deep representation learning, the stop-gradient operator ($\text{detach}(\cdot)$) has played an indispensable role in preventing representational collapse. Most prominently in Siamese self-supervised learning frameworks like SimSiam~\cite{chen2021exploring} and BYOL~\cite{grill2020bootstrap}, severing gradient flow along one branch prevents trivial constant solutions without requiring negative pairs. Similarly, in reinforcement learning and actor-critic methods, target networks and detached value estimators are crucial to prevent divergence in recursive temporal-difference updates~\cite{mnih2015human}. 

Despite its success in these fields, the principle of \textbf{gradient decoupling} has not previously been systematically applied to recurrent feedback segmentation networks, where prediction masks have historically been backpropagated without constraint. Our work formalizes the Feedback Firewall principle, demonstrating that gradient decoupling is essential for robust, generalizable recurrent medical vision architectures.
